\documentclass{article}
\usepackage{amsmath}
\usepackage{graphicx}
\usepackage{float} % Required for [H]
\usepackage{hyperref}
\usepackage[en-GB]{datetime2}
\DTMlangsetup[en-GB]{ord=omit}
\hypersetup{
    colorlinks=true,
    linkcolor=blue,
    urlcolor=blue
}

\title{Epidemiology - Week 5 - Measuring Disease Occurrence \& Effect}
\author{Fernaldy}
\date{\today}

\begin{document}

\maketitle

\section{Measurement of Occurrence}
\subsection{Incidence}
\textbf{Incidence Rate} is the number of \textbf{\textit{new}} cases of a disease occurring
in the \textbf{\textit{population at risk}} over a time period. Incidence expresses the risk of becoming ill and is useful for studying causation. Incidence starts with the population that is at risk but disease-free. Incidence is generally calculated by
\begin{equation*}
    \text{Incidence} = \dfrac{\text{Number of new events in a specified period}}{\text{Number of persons exposed to risk during this period}}
\end{equation*}
\subsubsection*{Cumulative incidence}
Cumulative incidence, like incidence but differs on the denominator where it instead takes the number of people at the beginning of the study.
\begin{equation*}
    \text{Incidence} = \dfrac{\text{Number of new events in a specified period}}{\substack{\text{Number of people free of the disease in the}\\\text{population at risk at the beginning of the period}}}
\end{equation*}

\subsection{Prevalence}
Prevalence indicates the total disease burden since it calculates both old and new cases. Prevalence estimates the probability of the population being ill at the period of study. It is useful to study how much burden the disease carries on the population at risk. Prevalence is generally calculated as follows,
\begin{equation*}
    \text{Prevalence} = \dfrac{\text{Number of people with the outcome at a specified time}}{\text{Number of people in the population at risk at the specified time}}
\end{equation*}
Furthermore, there are two types of prevalence rates:
\begin{itemize}
    \item \textbf{Point Prevalence Rate} is the number of cases of a disease at \textit{one point in time} (point)
    \item \textbf{Period Prevalence Rate} is the number of cases (both old and new) occurring at any point over a specified longer period in the \textbf{\textit{population at risk}}
\end{itemize}

\subsubsection*{Interrelationships of prevalence and incidence}
Prevalence is dependent on both incidence and disease duration. If prevalence is low and does not vary significantly with time, it can be calculated by:
\begin{equation*}
    \text{Prevalence} = \text{Incidence} \times \text{Average duration of disease}
\end{equation*}
A keyword that may be found on cases like this is usually \textit{\textbf{steady-state conditions}}, which means the rate at which new cases appear (incidence) is roughly equal to the rate at which cases disappear / finish (recover or mortality). This explains why long-term diseases can have high prevalence despite low
incidence.

\subsection{Population at Risk}
In a nutshell for incidence rate, population at risk just means it's the people that could potentially develop the disease / outcome of interest.
\\
The criterias are:
\begin{itemize}
    \item A member of the population (time, place, person)
    \item Without the disease of interest at the beginning of follow-up
    \item Biologically at risk of developing the disease of interest
\end{itemize}

\subsubsection*{Example Case: Hip Fracture}
Estimate for Scotland in 2024 for over 50s is 2,293,574. Incident cases were 8426. The incident rate would therefore be,
\begin{equation}
    \frac{8426}{2293574} \times 1000 = 3.7 \text{ cases per 1000 population at risk}
\end{equation}
\\
The formal formula for the rate is (Specify the time period, e.g. 1 year):
\begin{equation}
    Rates = \frac{Numerator}{Denominator}
\end{equation}
where
\begin{itemize}
    \item Numerator - Cases of interest
    \item Denominator - Population at risk
\end{itemize}
The rate itself is typically expressed per 1000 or 100,000.
\\
\subsection{Incidence Rate}
Incidence rate can then be calculated with the following formula:
\begin{equation}
    \text{Incidence Rate} = \frac{\text{New cases of disease (n)}}{\text{Population at risk}}
\end{equation}

\paragraph{There were 500 cases of disease X in region Y in 2024. The population of said region is 500,000 in 2024.
The annual incidence of disease X is 500. Calculate the incidence rate.}
\begin{equation*}
    \text{Incidence Rate} = \frac{500}{500000} = \text{1 case per 1000 population at risk}
\end{equation*}
Therefore, the incidence rate of disease X in region Y in 2024 is 1 case per 1000 population at risk.

\subsection{Prevalence Rate}
Since prevalence is a combination of all cases (old and new), the prevalence rate can then be calculated by
\begin{equation}
    \text{Prevalence Rate} = \frac{\text{All cases}}{Population}
\end{equation}


\subsubsection*{Example Case:}
On a particular day there were 5000 people already diagnosed with disease X living in region Y (Total population 500,000).
Calculate the prevalence of disease X on that day.
\begin{equation*}
    \text{Prevalence Rate} = \frac{5000}{500000} = 0.1 \text{ or 10 cases per 1000 population at risk}
\end{equation*}

\subsubsection{Point Prevalence}
Point prevalence is the measure of proportion of people in a population who have a disease at \textbf{a particular time}
(such as a particular date). Alternatively, it can be identified as a snapshot of the disease in time. The rate
can be calculated by
\begin{equation*}
    \textit{Point Prevalence} = \frac{\text{Number of existing cases on a specific date}}{\text{Number of people in the population at this date}}
\end{equation*}

\subsubsection{Period Prevalence}
Period prevalence on the other hand is the measure of proportion of people in a population with a disease \textbf{over a
specific period of time}. For example, the amount of people who had flu in winter 2006. The rate can be calculated by
\begin{equation*}
    \textit{Period Prevalence} = \frac{\text{Number of cases that occurred in a given period}}{\text{Number of people in the population during this period}}
\end{equation*}


\pagebreak
\section{Crude and Standardised Rates}
\subsection{Crude Rates}
The crude rate is typically expressed per 1000 or 10,000 of population and can be calculated by
\begin{equation}
    \text{Crude Rate} = \frac{\text{Number of events}}{\text{Total people in the population}}
\end{equation}
However, when measuring crude rates, population structure is typically not taken into account which could
blur the interpretation of said crude rates. But it is still useful for knowing the overall picture.

\subsubsection*{Crude mortality rate}
Crude mortality rate (or the death rate) for all deaths or a specific cause of death is calculated as follows:
\begin{equation*}
    \text{Crude mortality rate} = \dfrac{\text{Number of deaths during a specified period}}{\text{Number of persons at risk of dying during the same period}}
\end{equation*}

\subsubsection*{Example Case:}
Town A and B have both the population count of 20,000. The crude annual mortality rate is both 25 per 1000.
But 50\% of Town A are over 65 years old while only 20\% of Town B are over 65 years old. Both towns have
the crude annual mortality rate of 25 per 1000 but more specifically in Town B, younger people are dying
more than in Town A. The crude rate doesn't take into account the population structure (age distribution of
each towns' population).

\subsection{Standardised Rates}
Standardised Rates however does take into account the downside of Crude Rates by 2 possible methods:
\textbf{Direct Standardisation} or \textbf{Indirect Standardisation}

\subsubsection{Direct Standardisation}
Direct Standardisation compares the two by standardising the population. To do this, it introduces a
third population (or a reference population) such as \textit{The European Standard Population, ESP}.

\subsubsection*{Example:}
\begin{itemize}
    \item Town A: 40,000 of $<$65 y.o. and 60,000 of $\geq$65 y.o.; 1300 deaths per 100,000 people
    \item Town B: 80,000 of $<$65 y.o. and 20,000 of $\geq$65 y.o.; 1100 deaths per 100,000 people
\end{itemize}
From this information itself, we can see the mortality rate is higher in Town A than Town B. However,
Town A has a much higher proportion of elderly residents. Direct Standardisation can be used here to
determine which town has worse outcome after taking into account the age distribution.
\begin{table}[H]
\centering
\begin{tabular}{|l|l|l|l|}
\hline
\textbf{Age Group}   & \textbf{Town A Population} & \textbf{Town A Deaths} & \textbf{Town A Rate (per 1,000)} \\ \hline
\textbf{Young ($<$65)} & 40,000            & 80            & 2.0                     \\ \hline
\textbf{Old ($\geq$65)}  & 60,000            & 1220          & 20.33                   \\ \hline
\textbf{Total}       & 100,000           & 1300          & 13.0 \textbf{(Crude)}            \\ \hline
\end{tabular}
\caption{Town A}
\label{tab:town-a}
\end{table}

\begin{table}[H]
\centering
\begin{tabular}{|l|l|l|l|}
\hline
\textbf{Age Group}   & \textbf{Town B Population} & \textbf{Town B Deaths} & \textbf{Town B Rate (per 1,000)} \\ \hline
\textbf{Young ($<$65)} & 80,000            & 240            & 3.0                     \\ \hline
\textbf{Old ($\geq$65)}  & 20,000            & 410          & 20.5                   \\ \hline
\textbf{Total}       & 100,000           & 650          & 6.5 \textbf{(Crude)}            \\ \hline
\end{tabular}
\caption{Town B}
\label{tab:town-b}
\end{table}

Choosing the Standard (Reference) Population of 100,000 people with an even split 50\:50 of young and old.
We can calculate the deaths in both Town A and B as if it had the standard population by the following,
\\~\\
Town A
\begin{itemize}
    \item Young ($<$65): $50,000 \times \frac{2.0}{1,000} = 100 \text{ deaths}$
    \item Old ($\geq65$): $50,000 \times \frac{20.33}{1,000} = 1016.5 \text{ deaths}$
    \item Total expected deaths: $100 + 1016.5 = 1116.5 \text{ deaths}$
\end{itemize}

Town B
\begin{itemize}
    \item Young ($<$65): $50,000 \times \frac{3.0}{1,000} = 150 \text{ deaths}$
    \item Old ($\geq65$): $50,000 \times \frac{20.5}{1,000} = 1025 \text{ deaths}$
    \item Total expected deaths: $150 + 1025 = 1175 \text{ deaths}$
\end{itemize}

Then we can calculate the \textit{Age-Standardised Rate (ASR)} by dividing by the total standard population,
followed with converting to a per 1000 rate.
\begin{itemize}
    \item Town A: $\frac{1116.5}{100,000} \times 1000 = 11.17 \text{ per 1000}$
    \item Town B: $\frac{1175}{100,000} \times 1000 = 11.75 \text{ per 1000}$
\end{itemize}

We can then summarise the comparison of the two towns' crude rate and standardised rate
\begin{table}[H]
\centering
\begin{tabular}{|l|l|l|l|}
\hline
\textbf{Metric                     } & \textbf{Town A}         & \textbf{Town B}         \\ \hline
\textbf{Crude Rate                 } & 13.0 per 1000  & 11.17 per 1000 \\ \hline
\textbf{Age-Standardised Rate (ASR)} & 11.17 per 1000 & 11.75 per 1000 \\ \hline
\end{tabular}
\caption{Comparison of both town's crude and standardised rate}
\label{tab:comparison-metric}
\end{table}

The interpretation of both now becomes:
\begin{itemize}
    \item Crude Rate: Unadjusted and misleading due to differing population structure (age)
    \item ASR: Adjusted by age; Town B has slightly higher mortality risk
\end{itemize}

\subsubsection{Indirect Standardisation}
Indirect Standardised is instead used to compare the observed number of events within a population
of interest to those which would be expected if the standard population's event rates were applied
to the study population's age and sex structure.
\\~\\
When comparing mortality ratio, the \textbf{ Standardised Mortality Ratio (SMR) } is produced and the
SMR of a population is 100.
\begin{itemize}
    \item SMR $<$ 100 means there are fewer deaths than expected in the study population
    \item SMR $>$ 100 means there are more deaths than expected in the study population
\end{itemize}

\noindent This is typically used when
\begin{itemize}
    \item The age-specific event / death rates in the study population is unknown or unreliable; E.g. small sample, rare disease
    \item Only the total observed cases / deaths and the population size by age group is known for the study population
\end{itemize}
and can be calculated by
\begin{equation}
    \text{SMR} = \frac{\text{Observed Cases in Study Population}}{\text{Expected Cases if Study Population had Standard Rates}} \times 100
\end{equation}
Some conventions don't multiply by 100 and keep it at baseline 1.0 instead.
\\~\\
Where the expected cases is calculated by
\begin{equation}
\begin{split}
    \text{Expected Cases} = {} & \sum (\text{Study Population in Age Group } i \\
    & \times \text{Standard Population Rate in Age Group } i)
\end{split}
\end{equation}


\pagebreak
\section{Measures of Effect \& Difference}
What's typically measured here are two types: \textbf{Effect} or \textbf{Difference}.

\subsection{Effect}
\begin{itemize}
    \item Risk ratio (Relative Risk)
    \item Rate ratio
    \item Odds ratio
\end{itemize}

\subsubsection{Risk Ratio (Relative Risk)}
Risk ratio is the ratio of the probability of an outcome in an exposed group vs the probability of an outcome
in the unexposed group. Risk ratio is appropriate for fixed follow-up periods.

\subsubsection*{Example:}
The table below is a (hypothetical) prospective cohort study of lung cancer in people who smoke vs those
who don't smoke, followed up for 15 years.
\begin{table}[H]
\centering
\begin{tabular}{l|l|l|l|}
\cline{2-4}
                                 & Lung Cancer     & No Lung Cancer  & Risk           \\ \hline
\multicolumn{1}{|l|}{Smoker}     & 40 \textbf{ (a)} & 60 \textbf{ (b)} & 40\% (40 / 100) \\ \hline
\multicolumn{1}{|l|}{Non-Smoker} & 5 \textbf{(c)}  & 95 \textbf{(d)} & 5\% (5 / 100)   \\ \hline
\end{tabular}
\label{tab:lung-cancer}
\caption{Lung Cancer}
\end{table}
Where the smoker (exposed) or non-smoker (unexposed) are the groups and lung cancer status is the outcome of interest.
\\
The calculations for the data above is:
\begin{itemize}
    \item $\text{Risk for smokers} = \frac{40}{40 + 60} = 40\% $
    \item $\text{Risk for non-smokers} = \frac{5}{5 + 95} = 5\% $
    \item $\text{Risk Ratio} = \frac{40\%}{5\%} = 8 \text{ over 15 years}$
\end{itemize}
In other words: \textit{Over a 15 year period, smokers are 8 times more likely to develop lung cancer
compared to non-smokers.}


\subsubsection{Odds Ratio}
Odds ratio is the ratio of the probability of an event occurring vs the probability of it not occurring.
Using the table from the \textit{Risk Ratio} section, the odds ratio can be calculated by:
\begin{itemize}
    \item $\text{Odds of smoker developing lung cancer:} = \frac{40}{60} = 0.67$
    \item $\text{Odds of non-smoker developing lung cancer:} = \frac{5}{95} = 0.05$
    \item $\text{Odds Ratio} = \frac{40 \times 95}{60 \times 5} = 12.67$
\end{itemize}
In other words: \textit{The odds of developing lung cancer among smokers is 12.67 times the odds of
developing lung cancer among the non-smokers over the 15 year period.}
\\
However, odds ratio tend to over-estimate risk. Furthermore, odds ratio can be used when the
denominator is not known, such as in a case-control study. When there is a low number of events,
the odds ratio tend to be similar to the risk ratio.


\subsubsection{Rate Ratio}
Rate ratio is the comparison of the incidence rate in one group vs the incidence rate in another group
of the same event. The measure differs from risk ratio by its denominator, instead of the individuals at risk
it uses the total-person time at risk. The rate ratio is typically used when there's a time aspect in the study,
such as when participants are followed-up at different times.



\subsection{Difference}
\begin{itemize}
    \item Risk difference
    \item Rate difference
\end{itemize}

\subsubsection{Risk Difference}
Risk difference (Attributable risk or excess risk) is the absolute difference in risk of an outcome
between the exposed an unexposed group. Using the same \hyperref[tab:lung-cancer]{lung cancer table},
the risk difference can be calculated by subtracting the risk of lung cancer in non-smokers from the
risk of lung cancer in smokers.
\begin{equation*}
    \text{Risk Difference} = 40\% - 5\% = 35\%
\end{equation*}

The purpose of risk difference is to quantify the amount of risk that can be attributed to the exposure.
In other words: \textit{Smoking leads to an absolute increase of 35\% in the risk of developing
lung cancer over a 15 year period.}


\subsubsection{Rate Difference}
Rate difference compares two incidence rates. Take the following example:

\paragraph{If an exposed group had an incidence rate of 40 new cases per 1,000 person-years and the
unexposed group had an incidence rate of 15 new cases per 1,000 person-years. Calculate the rate difference.}
\begin{equation*}
    \text{Rate Difference} = 40 - 15 = 25 \text{ new cases per 1,000 person-years}
\end{equation*}
The purpose of rate difference is to calculate the burden of disease attributable to factor. What this means is
given the example above, 15 cases would have happened anyway (baseline risk). However, the remaining 25 cases
can be directly attributed to the exposure.

\subsection{Risk Factors and Protective Factors}
Risk factor is any condition, behaviour, or characteristic that increases the probability of a particular disease
or outcome occurring. This factor can be biological, behavioural, environmental, or genetic.

Protective factor is the opposite of risk factor where instead of increasing the probability, it reduces the
disease or outcome from occurring.


\subsection{Statistical Association}
Statistical association is the relationship / link between two or more variables where their patterns of variation
are related. For example smoking is related to lung cancer, when the smoking prevalence in a population increases
the incidence of lung cancer also increases.

\pagebreak
\section{Examples on Calculating Incidence and Prevalence}
\subsection{Example 1}
A district public health team wants to understand the burden of chronic obstructive
pulmonary disease (COPD) in their community. They conduct a door-to-door health
survey of 2,400 adults during March 2025. At the time of the survey, 192 respondents
report a current clinical diagnosis of COPD (confirmed by their medical records).
\\
Calculate the point prevalence of COPD in this adult population.
\\
\begin{equation*}
    \text{Prevalence} = \dfrac{192}{2400} = 0.08 \text{ or 8\%}
\end{equation*}
Interpretation: As of this time (March 2025), 8\% of the population has COPD.

\subsection{Example 2}
In 2024, an outbreak of norovirus occurred across three care homes. Researchers
followed 600 residents from January to December to determine the risk of infection. By
the end of the year, 90 residents had experienced at least one confirmed episode of
norovirus.
\\
Calculate the cumulative incidence
\\
\begin{equation*}
    \text{Cumulative incidence} = \dfrac{90}{600} = 0.15 \text{ or 15\%}
\end{equation*}
Interpretation: Over the course of 2024, the cumulative incidence (risk) of confirmed norovirus infection among care home residents was 15\%.

\subsection{Example 3}
A 2-year prospective study of 400 haemodialysis patients investigates the occurrence
of bloodstream infections. During follow-up, some patients transfer or die, resulting in a
total of 650 person-years of observation. There are 52 new bloodstream infections
documented.
\\
Calculate the incidence rate per 1,000 person-years.
\\
\begin{equation*}
    \text{Incidence rate} = \dfrac{52}{650} = 0.08 \text{ or 80 per 1,000 person-years}
\end{equation*}

\subsection{Example 4}
A chronic neurological condition, multiple sclerosis (MS), affects adults in a stable
regional population of 500,000 people. Epidemiologists estimate that 100 new cases
occur each year. The average duration from diagnosis to death is 20 years.
\\
Estimate the prevalence of MS in this population under steady-state conditions (consistent per year).
\\
Since it is under steady-state conditions,
\begin{equation*}
    \text{Prevalence} = \dfrac{100}{500,000} \times 20 = 0.04 \text{ or 4\%}
\end{equation*}

\subsection{Example 5}
A national initiative offers free blood pressure checks in pharmacies. Over 12 months,
100,000 adults are screened. Among them, 8,000 are newly diagnosed with
hypertension. However, many people already knew they had hypertension before the
campaign began, an additional 12,000 participants report being on antihypertensive
medication at baseline.
\\
Should we describe the finding as an incidence or prevalence estimate? What would
the question look like for either of these?
\\
Both can be calculated, it would depend on the RQ
\begin{itemize}
    \item If the question is \textit{How many people currently have hypertension?}
        \begin{equation*}
            \text{Prevalence} = \dfrac{8000 + 12000}{100000} = 0.2 \text{ or 2\%}
        \end{equation*}
    \item If the question is \textit{How many new cases were detected during the screening?}
        \begin{equation*}
            \text{Incidence} = \dfrac{8000}{100000 - 12000} = 0.0909 \text{ or 9.1\%}
        \end{equation*}
\end{itemize}

\subsection{Example 6}
A mental health researcher surveys 2,000 first-year university students in October. At
baseline, 300 students report symptoms consistent with major depression (PHQ-9 $\geq$ 10).
A follow-up survey six months later finds 160 new cases among those who were
symptom-free initially, while 200 of the original still meet criteria.
\\
Which measure(s) should be reported — prevalence, incidence, or both? Calculate the relevant values.
\\
Both can be calculated, it would depend on the RQ
\begin{itemize}
    \item Baseline
        \begin{equation*}
            \text{Prevalence} = \dfrac{300}{2000} = 0.15 \text{ or 15\%}
        \end{equation*}
    \item Cumulative incidence (6 months)
        \begin{equation*}
            \text{Cumulative incidence} = \dfrac{160}{2000 - 300} = 0.094 \text{ or 9.4\%}
        \end{equation*}
    \item Point prevalence (6 months)
        \begin{equation*}
            \text{Point prevalence} = \dfrac{160 + 200}{2000} = 0.18 \text{ or 18\%}
        \end{equation*}

\end{itemize}

\end{document}
